\section{Conclusion}
\label{sec:conclusion}

We presented a strong applied result for stop-level transit demand
classification on the ZTBus dataset~\cite{widmer2023ztbus}: a Random
Forest trained with stop-level lag features reaches macro-F1 = 0.8287
on the realistic temporal evaluation (pre-2022 train, 2022 test), a
$+0.2221$ improvement over a 57-feature cross-sectional baseline of
GPS, speed, power, and calendar features. The three stop-level lag
features (\texttt{pax\_stop\_lag\_1}, \texttt{pax\_stop\_lag\_3},
\texttt{pax\_stop\_lag\_5}) contribute approximately 22\% of total
Gini importance in the winning Random Forest, with lag~1 dominating
as predicted by the PACF analysis.

Along the way we documented an invisible target leakage pattern that
arises from the interaction of two standard time series preprocessing
steps, forward-fill imputation and lag feature computation, when applied
to sparse event data. A naive 1~Hz-lag-on-forward-filled-data approach
inflates Decision Tree pooled macro-F1 from 0.5713 (baseline) to
0.8743 ($+0.30$), apparently indicating strong temporal predictive
power;
diagnosis reveals the lag features equal the current target value for
approximately 95\% of rows due to the forward-fill staircase, making
them near-copies of the target variable. The leakage survives all
standard safeguards (mission-level splits, honest cross-validation, no
test-set contamination) because it is endogenous to the feature
construction on the training data alone. Our fix computes lag features
at the natural stop-event granularity and forward-fills the \emph{lag
values} onto the dense grid rather than forward-filling the raw series
first, eliminating the leakage while preserving the physically
meaningful autoregressive signal.

We further tested whether TimesFM~2.5-200M~\cite{das2024timesfm}
penultimate-layer embeddings could improve over stop-level lag features
and obtained a null result: no model benefits from the embeddings. The
100\% PCA explained variance at 32 components suggests the foundation
model's representations for this domain live in a low-rank subspace that
is already spanned by short-range stop history. This is a cautionary
counterweight to the assumption that foundation models always help when
used as feature extractors on tabular time series data.

Future work includes (1) testing the leakage pattern and fix on other
sparse event domains (healthcare vital signs, financial tick data, IoT
sensor streams), (2) fine-tuning TimesFM on the ZTBus data to see if
task-adapted embeddings carry information beyond hand-crafted lag
features, (3) deploying the model in a real-time transit operations
setting with an explicit future-horizon forecast, (4) extending to a
regression formulation to better model the high-demand tail, and (5)
estimating performance variance across multiple random seeds (the
present paper uses a single fixed training seed). Our reproducible
pipeline, released as open source, provides a starting point for each
of these directions.
